\section{Introduction}\label{sec:intro}

The prime factors of a typical integer, viewed on a logarithmic
scale, form a random partition of unit mass.  We prove that the same
limiting partition occurs for the predecessor of a uniformly chosen
prime.

For a prime $p\ge3$, list the prime factors of $p-1$, with
multiplicity, in decreasing order,
\[
 q_1(p)\ge q_2(p)\ge\cdots,
\]
and put $q_j(p)=1$ after the list is exhausted.  Define
\[
 V_j(p)=\frac{\log q_j(p)}{\log(p-1)}.
\]
Then $V_j(p)\ge0$ and $\sum_jV_j(p)=1$.

Let $U_1,U_2,\ldots$ be independent uniform random variables on
$(0,1)$, and form the stick fragments
\begin{equation}\label{eq:stick-breaking}
 B_1=1-U_1,
 \qquad B_j=\left(\prod_{i<j}U_i\right)(1-U_j)
       \quad(j\ge2).
\end{equation}
Their decreasing rearrangement $(L_1,L_2,\ldots)$ has the
Poisson--Dirichlet distribution with parameter one, denoted
$\PD(1)$. Write $\pi(x)$ for the number of primes at most $x$.

\begin{theorem}\label{thm:main}
For every fixed $k\ge1$ and every bounded continuous function
$F:[0,1]^k\to\R$,
\begin{equation}\label{eq:main-law}
 \lim_{x\to\infty}\frac1{\pi(x)-1}
     \sum_{\substack{3\le p\le x\\p\text{ prime}}}
       F\bigl(V_1(p),\ldots,V_k(p)\bigr)
 =\E F(L_1,\ldots,L_k).
\end{equation}
The limit holds through all real $x$ and uses ordinary equal
weighting of the primes.
\end{theorem}

Theorem~\ref{thm:main} resolves positively the conjecture of Ford,
Konyagin and Luca \cite[Section~6, Conjecture~5]{FordKonyaginLuca2010}.
Their formulation uses the same multiplicity convention, normalization,
and ordinary counting distribution on primes. The assertion includes
all finite joint distributions, strengthening the largest-factor
prediction alone.

\paragraph{History and significance.}
For ordinary uniformly sampled integers, the largest-factor law begins
with Dickman \cite{Dickman} and the smooth-number estimates of
de Bruijn \cite{deBruijn1951}. Billingsley proved the joint limiting
law of the ordered large prime divisors \cite{Billingsley}.
Donnelly and Grimmett \cite[Theorem~1 and Corollaries~2--3]{DonnellyGrimmett}
gave a size-biased treatment, including multiplicities, leading directly
to the stick-breaking and Poisson--Dirichlet descriptions. Arratia,
Kochman and Miller \cite{ArratiaKochmanMiller2014} express this
identification through the joint intensities of distinct factor
tuples. Our final passage from interior factor statistics to ranked
factors uses the same size-biased approach; that passage is proved
explicitly in Section~\ref{pd:section}.

The shifted-prime problem has a different arithmetic difficulty:
requiring a product of large primes to divide $p-1$ places $p$ in
a progression with that product as modulus. Standard distribution
estimates give useful information for products below a fixed power
of $x$, whereas the full factor partition involves products throughout
the range below $x$. The smoothness question arose in Erd\H{o}s's
work on Euler's function \cite{Erdos1935,Erdos1956}, and
Pomerance \cite{Pomerance1980} developed its connection with large
totient fibers. Baker and Harman \cite[Theorem~1]{BakerHarman1998}
and Lichtman \cite[Theorem~1.1]{Lichtman2022} obtained lower bounds
of the form $x/(\log x)^C$ for primes with smooth predecessors, at
thresholds $x^{0.2961}$ for $p\le x$ and $x^{0.2844}$ for
$x<p\le2x$, respectively.

For shifted primes, Granville states the conjectural asymptotic
\cite[Section~5.3, Equation~(5.1)]{Granville}
\begin{equation}\label{eq:shifted-dickman}
 \#\{p\le x:P^+(p-1)\le x^{1/u}\}
       \sim\pi(x)\rho(u)\qquad(u\ge1\text{ fixed}),
\end{equation}
where $P^+(n)$ is the largest prime factor, with $P^+(1)=1$, and the Dickman function
is determined by $\rho(u)=1$ for $0\le u\le1$ and
$u\rho'(u)+\rho(u-1)=0$ for $u>1$.
Theorem~\ref{thm:main} resolves this fixed-$u$ conjecture
positively and proves the complete joint limiting law.  The use of
$x$ rather than $p-1$ in the threshold in
\eqref{eq:shifted-dickman} causes no change: restrict first to
$\varepsilon x<p\le x$, where
$\log(p-1)/\log x\to1$ uniformly, and then let
$\varepsilon\downarrow0$.  The largest part of $\PD(1)$ has the
continuous Dickman distribution, by the ordinary-integer limit in
\cite[Section~1]{DonnellyGrimmett} and the smooth-number asymptotic
in \cite[Section~1, Equation~(1.1)]{Granville}; hence the
corresponding threshold sandwich applies.

The closest general distribution theorem is due to Bharadwaj and
Rodgers \cite[Theorem~7]{BharadwajRodgers2026}. Their
Poisson--Dirichlet theorem applies to nonnegative sequences satisfying
their regularity and congruence-uniformity hypotheses, with level of
distribution one. For shifted primes this gives the
full law under the Elliott--Halberstam conjecture, as anticipated in
\cite[Section~6]{FordKonyaginLuca2010}. Unconditionally,
\cite[Proposition~2 and Lemma~8]{BharadwajRodgers2026} give level
one half and the corresponding factor correlations for tests supported
where the sum of the logarithmic coordinates is less than one half.
Our extraction reaches every compact subset of the open simplex with
coordinate sum less than one, and the final probability argument
controls its boundary. All smoothness parameters in
\eqref{eq:shifted-dickman} remain fixed as $x$ grows.

There is also strong recent information about the small prime
divisors of shifted primes. Ford \cite{Ford2025} proves a
total-variation approximation for their valuations on subpower ranges;
Gorodetsky \cite{Gorodetsky2026} obtains such approximations within
a general sieve model. The large-factor law studied here concerns
the logarithmic mass carried by prime factors of polynomial size.

Ford, Konyagin and Luca formulated their conjecture in studying
prime chains and Pratt trees. Theorem~\ref{thm:main} supplies the
factorization law at the root of their probabilistic model; the
recursive independence assumptions used for later generations are
additional hypotheses. Smooth predecessors also underlie the
Carmichael-number construction of Alford, Granville and Pomerance
\cite{AlfordGranvillePomerance1994}, together with a separate
progression-distribution input.

The companion article \emph{Weighted dilation graphs, smooth shifted
primes and totient fibers}~\cite[Theorem~1.2]{SmoothShifted}, denoted S, proves that, for every
fixed positive smoothness exponent, there are $x^{1-o(1)}$ primes
in its stated interval with smooth predecessors.  Such a lower
bound does not imply a positive limiting proportion, much less
\eqref{eq:main-law}. We use S's dilation-graph operator theorems in a
different extraction argument. Section~\ref{sec:correlations} states
their precise contracts and verifies the new endpoint hypotheses.
The long-prime polynomial and logarithmic-phase estimates needed in
both analytic branches are proved in Appendix~\ref{lp:section};
Section~\ref{sec:inputs} records the ranges in which we use them.

\paragraph{The new analytic step.}
The main new estimate is Theorem~\ref{thm:type-II}: a marked
Type~II estimate in which only the predecessor side carries
divisor marks, namely weights recording selected small prime divisors.
It permits a prime factor on the other side to be
replaced by a rough integer, free of primes below a prescribed cutoff,
with its local density. Cauchy's
inequality, applied after dividing out one tuple of marks, leads
to a small-determinant relation between primitive lattice vectors.
We average its long signed moments using a root-residue formula
and an exact memory expansion.  The memory records a prime across
gaps between its uses, so its divisibility probability is charged
only once.

The root coordinates are restricted by quantitative separation
conditions. These give both a balanced lattice
box for fresh-label minor arcs and a separation rule for retrieving
remembered labels. Global distinctness is then restored by grouping
equalities between newly drawn prime labels according to their rank,
the number of independent equality constraints. High ranks supply
many small point-probability factors, whereas low ranks
alter only a vanishing fraction of the signed edge contractions.
The resulting estimate has arbitrarily strong fixed logarithmic
precision, with a number of marking bands independent of their
particular fixed exponents.

The signed-moment method has precedents in divisibility graphs.
Helfgott and Radziwi\l\l{} center prime-divisibility weights by
subtracting $1/p$ and analyze high traces through prime-labelled
walks \cite[Sections~1.3, 1.5 and~5]{HelfgottRadziwill2021}.
Their proposed composite shifts \cite[Section~9.1]{HelfgottRadziwill2021}
are developed by Pilatte using products of primes and a
non-backtracking trace estimate
\cite[Definition~3.2 and Proposition~5.3]{Pilatte2026}.
Recurring prime labels create dependencies in these walk expansions.
Here the state contains active prime lists and a primitive lattice
vector. At each prime $p$, root averaging gives a uniform projective
line, so a specified line has probability $1/(p+1)$. The memory
preserves that line across gaps between active runs.
Sections~\ref{sec:determinant}--\ref{mem:section}
develop this determinant geometry and prove the exact memory identity;
Section~\ref{sec:major-arcs} supplies the complementary major-term
estimate.

\paragraph{From marked correlations to the full law.}
The second estimate, Theorem~\ref{thm:two-sided}, uses marks on
both endpoints.  We apply the general graph inputs of S to a
centered long-prime divisor statistic and prove the required new
endpoint Fourier estimates.  A small mark restricts difficult
frequencies to a sparse set; a long-prime polynomial controls the
positive tuple term there, and a finite divisor model controls
the constant term.

A first-moment presieve and the one-sided estimate replace all
prime slots of the composite terms by rough slots.  Successful
marking groups then bring these composite terms within the
two-sided estimate.  Averaging over disjoint fixed marking arrays
removes the marks from the prime average.  This proves interior
factorial-moment asymptotics for ordinary prime dyads.

Finally, sequential size-biased sampling turns the interior
factorial measures into probability densities of total mass one.
This excludes escape to the boundary and yields the stick fragments
in \eqref{eq:stick-breaking}.  The expected undrawn mass controls
sorting, and a finite dyadic decomposition gives all real upper
endpoints.  We provide this probability argument explicitly, since
restricted-support moment formulas alone are not the statement of
Theorem~\ref{thm:main}.
Figure~\ref{fig:proof-map} summarizes the dependence of these steps.
The two-sided endpoint analysis is in Section~\ref{sec:correlations},
the prime extraction and removal of marks in Section~\ref{ext:section},
and the probability argument in Section~\ref{pd:section}.
\begin{figure}[tb]
\centering
\begin{tikzpicture}[
  box/.style={draw=black!65,rounded corners=2pt,align=center,
    text width=5.55cm,minimum height=1.03cm,inner sep=5pt,font=\small},
  arr/.style={-{Latex[length=2mm]},thick},
  node distance=8mm and 7mm]
\node[box,fill=blue!5] (det)
  {Determinant graph and signed memory\\
   Sections~\ref{sec:determinant}--\ref{mem:section}};
\node[box,fill=blue!5,right=of det] (dil)
  {Dilation-graph inputs and endpoint estimates\\
   Section~\ref{sec:correlations}};
\node[box,below=of det] (one)
  {One-sided marked correlation\\Theorem~\ref{thm:type-II}};
\node[box,below=of dil] (two)
  {Two-sided marked correlation\\Theorem~\ref{thm:two-sided}};
\node[box,below=11mm of $(one.south)!.5!(two.south)$,
      text width=9cm] (interior)
  {Presieve, replacement of prime slots, and averaging of marks\\
   Interior factorial moments: Theorem~\ref{thm:interior}};
\node[box,below=of interior,text width=9cm,fill=blue!5] (pd)
  {Size-biased sampling, full mass, and decreasing rearrangement\\
   Poisson--Dirichlet limit: Theorem~\ref{thm:main}};
\draw[arr] (det) -- node[right,font=\scriptsize,align=left]
  {major-term estimate\\Section~\ref{sec:major-arcs}} (one);
\draw[arr] (dil) -- (two);
\draw[arr] (one.south) -- (interior.north west);
\draw[arr] (two.south) -- (interior.north east);
\draw[arr] (interior) -- (pd);
\end{tikzpicture}
\caption{The proof has two analytic inputs. The one-sided estimate
replaces prime slots in composite terms by rough slots. Successful
marking brings the resulting centered statistic under the two-sided
estimate. Averaging the fixed marks then gives interior moments;
the probability argument recovers the full law.
Appendix~\ref{lp:section} supplies the long-prime and
logarithmic-phase estimates used in the analytic branches.}
\label{fig:proof-map}
\end{figure}

